\chapter{QICAS active-space optimization (menu 24)}
\label{ch:menu24}
\menulabel{24}

\section{Purpose}

QICAS (Ding, Knecht and Schilling, JPCL 2023) optimizes the orbital basis of a
CAS($n$, $m$) scheme \emph{without touching the Hamiltonian}.  Its cost
function is the out-of-CAS correlation
$F_{\mathrm{QI}}(\mathcal{B}) = \sum_{i \notin A} S(\rho_i)$, the sum of the
four-state entropies of the non-active orbitals, and the source's Theorem~1
bounds the CASCI error by it:
$E_{\mathrm{CASCI}} - E_{\mathrm{FCI}} \le \Delta E_{\max}/\ln 4 \cdot
F_{\mathrm{QI}}(\mathcal{B})$.  $F_{\mathrm{QI}}$ depends only on the one- and
two-particle RDMs, so the orbitals can be rotated to minimize it cheaply, and
the CASCI in the optimized basis approaches the FCI.

\section{What you need}

The FCIDUMP of a converged CASSCF run (the same dump menu~\menu{12} reads),
the target active space \code{nel,norb} within the window, and --- optionally
--- the run's output, so the CI root is matched to the printed CASSCF energy
(the menu~\menu{12} cross-check).

\section{Walkthrough}

\lstinputlisting[style=fbkcode, firstline=29, lastline=61,
  title={The menu-24 example session on the N2 CAS(6,6) FCIDUMP: a deliberately
  awkward $(2e, 4o)$ request that splits the $\pi$ pair --- the Jacobi sweep
  finds the $F_{\mathrm{QI}}$ fixed point ($0.2481 \to 0.0302$), the occupancy
  reading reclassifies the space to the $(4e, 4o)$ $\pi$ arrangement, and the
  CASCI gap drops from $6.70\times10^{-2}$~Eh to $5.33\times10^{-3}$~Eh}]
  {examples/expected/m24_qicas.txt}

\section{What you get}

The per-orbital entropy profile before and after the optimization, the
$F_{\mathrm{QI}}$ values, the CASCI energies in the initial and the optimized
basis against the window FCI, the Theorem-1 check with the run's own numbers,
and a report (\file{<FCIDUMP>.qicas.fbk.md}) with the citations.

\section{How to read it}

\begin{readbox}
\begin{itemize}
  \item the partition is read two ways: the \emph{requested} one (this tool
    orders the dumped orbitals by natural occupation and cuts the highest as
    closed, the lowest as virtual) and the \emph{final} one --- the occupancy
    reading of the optimized basis (an orbital counts as closed if its
    occupancy ends above 1).  The optimizer may rotate a strongly correlated
    orbital out of a closed slot and an uncorrelated one in, so the final
    space can read as a different $(n, m)$: in the example the awkward
    $(2e, 4o)$ request is repaired to the $(4e, 4o)$ $\pi$ space and the CASCI
    gap improves by a factor of twelve;
  \item \textbf{the rotation set} decides how far the sweep may reach:
    \code{touch} (the default) rotates every pair touching a non-active
    orbital --- the source's chemical-accuracy choice; \code{exclusive}
    restricts to active/non-active pairs, its economical variant for large
    spaces.  $F_{\mathrm{QI}}$ is \emph{not} invariant under rotations within
    the non-active space, which is why the default includes them;
  \item the Theorem-1 line is a genuine check: when a target sits so far from
    the window ground state that the reference-state overlap deficit exceeds
    $1/2$ (the assumption of the source's proof), the report says the
    inequality is not informative instead of hiding it;
  \item \textbf{scope, stated in the report}: this is QICAS applied \emph{on
    the FCIDUMP window} (the source's own subset application) with exact RDMs
    from the four-state-entropy route.  The source drives it from a
    full-space DMRG ground state, which this toolbox does not have, so
    $F_{\mathrm{QI}}$ here is not comparable with the source's full-space
    values;
  \item the optimized rotation is reported as a matrix and used for the CASCI
    check; writing it back into a \code{.gbw} (the mkl route of
    menu~\menu{18}) is registered as a follow-up, as is the source's
    size-selection variant (its Appendix~C: minimize the total orbital
    entropy, read the plateau of the threshold diagram).
\end{itemize}
\end{readbox}
